\documentclass[11pt]{article}
\usepackage[margin=1in]{geometry}
\usepackage{amsmath,amssymb,amsthm}
\usepackage{graphicx}
\usepackage{booktabs}
\usepackage{hyperref}
\usepackage{xcolor}
\usepackage{natbib}
\usepackage{subcaption}

\newcommand{\citeneeded}[1]{\textcolor{red}{[CITE: #1]}}
\newcommand{\todo}[1]{\textcolor{orange}{[TODO: #1]}}
\newcommand{\numplaceholder}[1]{\textcolor{blue}{\textbf{#1}}}

\title{Exact Ground Truth for Epistasis Detection:\\A Benchmark from Combinatorially Complete Circuit Data}

\author{
  Elliot Tower \\
  \texttt{elliot@elliottower.ai}
}

\begin{document}
\maketitle

\begin{abstract}
Higher-order interactions (epistasis) drive phenotype in genetics and
function in neural circuits, yet every method for detecting them is
validated against other approximations---never against exact ground truth,
because exhaustive combinatorial measurement is infeasible at biological
scale. We construct an exact benchmark by computing all $2^{15} = 32{,}768$
coalition values over attention-head subsets in GPT-2-small, yielding
the complete Walsh--Hadamard interaction spectrum for three circuits
spanning near-additive, interactive, and null regimes. Against this
ground truth, we evaluate five interaction-detection methods from the
genetics and explainable-ML literatures---LASSO-Walsh, iterative random
forests, and three Shapley interaction approximators---across seven
sample budgets with pre-registered hypotheses.
%
The dominant factor separating methods is \emph{basis alignment}: whether
the method's representational assumptions match the underlying
interaction structure. LASSO-Walsh, which estimates Walsh coefficients
directly, achieves near-perfect pairwise interaction detection
(\numplaceholder{AUROC $> 0.99$}) at 5\% sample cost, while methods
operating through an intermediate basis (k-Shapley Interaction Indices)
reach \numplaceholder{$\sim$0.85} despite comparable reconstruction
accuracy. Three metric-design bugs caught during development---each a
named failure mode recurring across the interaction-detection
literature---are reported as methodological contributions. The benchmark,
pre-registration artifacts, and all coalition tables are released for
extension.
\end{abstract}


%% ====================================================================
\section{Introduction}
\label{sec:intro}
%% ====================================================================

% GPS: Goal -> Problem -> Solution

% GOAL: Detecting higher-order interactions (epistasis) is central to
% genetics and mechanistic interpretability.

Epistasis---the dependence of a phenotype on combinations of genetic
variants rather than individual effects---is among the oldest unsolved
measurement problems in quantitative genetics \citeneeded{Fisher 1918,
Bateson 1909}. Analogous higher-order interactions appear in neural
circuit analysis, where ablating individual components understates the
joint effect of removing a functionally coupled group
\citeneeded{Wang et al 2022 IOI, Conmy et al 2023 ACDC}.

% PROBLEM: No ground truth exists for method validation.

Every method for detecting epistasis---BOOST \citeneeded{Wan et al 2010},
MDR \citeneeded{Ritchie et al 2001}, Walsh-analysis approaches
\citeneeded{Neidhart et al 2013}, and Shapley interaction indices
\citeneeded{Grabisch Roubens 1999, Sundararajan et al 2020}---is validated
against other methods or against partial ground truth (pairwise only,
simulated, or measured on a few loci). The combinatorial explosion of
$2^n$ subset evaluations makes exhaustive measurement infeasible at
biological scale, so the field validates approximations against
approximations.

% SOLUTION: We construct exact ground truth at n=15.

Transformer circuits offer a setting where exhaustive evaluation is
tractable: $n = 15$ attention heads yield $2^{15} = 32{,}768$ coalitions,
each evaluable in a single forward pass over a prompt batch. We compute
all coalition values for three circuits in GPT-2-small, producing the
complete Walsh--Hadamard interaction spectrum---exact to floating-point
precision---for a near-additive circuit, an interactive circuit, and a
null control. Against this ground truth, we evaluate five
interaction-detection methods spanning the genetics and explainable-ML
lineages, with pre-registered hypotheses and falsification criteria.

% Close: state the thesis

The central finding is that \emph{basis alignment}---whether a method's
native representation matches the interaction structure being
scored---explains more of the performance gap between methods than
algorithmic sophistication or sample budget. This result has direct
implications for genetics practitioners choosing among interaction
methods: the choice of representational basis constrains recovery
fidelity before any algorithm runs.


%% ====================================================================
\section{Exact Ground Truth}
\label{sec:ground-truth}
%% ====================================================================

% GPS: describe what "exact" means and characterize the three regimes.

\subsection{Coalition value functions over circuit subsets}

\todo{Define the value function v(S) = logit\_diff under zero-ablation
of head subset S. Describe the IOI task. Cite Wang et al.}

For a set of $n = 15$ attention heads, we evaluate the value function
$v(S)$ for every $S \subseteq \{1, \ldots, n\}$ by ablating the
complement $\bar{S}$ and measuring the logit difference between the
correct and incorrect indirect-object tokens, averaged over 512 prompts.
This produces a vector $v \in \mathbb{R}^{2^{15}}$ from which the
complete Walsh--Hadamard spectrum is computed via the butterfly algorithm
in $O(n \cdot 2^n)$ time.

\subsection{Three interaction regimes}

\begin{table}[t]
\centering
\caption{Ground-truth characterization of the three benchmark circuits.
$k_{99}$ is the number of Walsh coefficients capturing 99\% of
non-intercept energy. The energy spectrum shows the fraction of total
Walsh energy at each interaction order. The three circuits span distinct
interaction regimes: near-additive (canonical IOI), interactive
(weight-identified), and null (random heads).}
\label{tab:circuits}
\begin{tabular}{lrrrrrr}
\toprule
Circuit & $n$ & $k_{99}$ & Order 0 & Order 1 & Order 2 & Order 3+ \\
\midrule
\texttt{weight\_ioi} & 15 & \numplaceholder{XX} & 97.5\% & 1.1\% & 0.7\% & 0.7\% \\
\texttt{canonical\_ioi} & 15 & \numplaceholder{XX} & 94.5\% & 3.5\% & 1.5\% & 0.5\% \\
\texttt{random15} & 15 & \numplaceholder{XX} & 99.3\% & 0.4\% & 0.2\% & 0.1\% \\
\bottomrule
\end{tabular}
\end{table}

\textbf{Weight-identified circuit (\texttt{weight\_ioi}).}
\todo{Describe provenance: identified via factorized weight decomposition,
not activation patching. Cite the factorization paper.}
This circuit has the richest higher-order structure among the three,
with \numplaceholder{XX\%} of non-intercept energy above order 2.
The $k$-SII reconstruction ceiling at order $\leq 2$ is $R^2 = 0.757$,
rising to $0.868$ at order $\leq 3$ (Table~\ref{tab:ceilings}).

\textbf{Canonical IOI circuit (\texttt{canonical\_ioi}).}
The 15 heads identified by \citeneeded{Wang et al 2022} via activation
patching. Near-additive: most non-intercept energy is at orders 1--2,
with a $k$-SII ceiling of $R^2 = 0.955$ at order $\leq 2$.

\textbf{Random control (\texttt{random15}).}
Fifteen heads drawn uniformly at random from all 144 GPT-2-small heads.
Serves as a negative control---any structure detected here is noise.
Near-perfect $k$-SII ceiling ($R^2 = 0.995$ at order $\leq 2$) reflects
the absence of genuine interaction structure.

\subsection{Truncation ceilings}
\label{sec:ceilings}

\begin{table}[t]
\centering
\caption{Exact $k$-SII reconstruction ceilings computed via
\texttt{shapiq.ExactComputer} v1.6.0. These are upper bounds on $R^2$
achievable by any method operating at the stated truncation order in the
$k$-SII basis. The gap between order-2 and order-3 ceilings quantifies
how much recoverable structure lies above pairwise interactions.}
\label{tab:ceilings}
\begin{tabular}{lcc}
\toprule
Circuit & $k$-SII $\leq 2$ & $k$-SII $\leq 3$ \\
\midrule
\texttt{weight\_ioi} & 0.757 & 0.868 \\
\texttt{canonical\_ioi} & 0.955 & 0.988 \\
\texttt{random15} & 0.995 & 1.000 \\
\bottomrule
\end{tabular}
\end{table}

The ceiling table is a first-class result: it quantifies the truncation
bottleneck that limits all order-bounded methods regardless of estimation
quality. On the interactive circuit, 24\% of reconstructable variance
is inaccessible to any order-2 method---a structural limitation no
algorithm can overcome.


%% ====================================================================
\section{Methods Under Test}
\label{sec:methods}
%% ====================================================================

% One paragraph per method: what basis/inductive bias it assumes.

\textbf{LASSO-Walsh.}
$L_1$-regularized regression on the Walsh basis functions up to order 4.
The design matrix $\Phi$ has columns $\phi_T(S) = (-1)^{|S \cap T|}$
for each Walsh index $T$ with $|T| \leq 4$. The LASSO coefficients are
direct estimates of the normalized Walsh coefficients
$\hat{w}_T = \frac{1}{2^n} \sum_S v(S)(-1)^{|S \cap T|}$. Regularization
path selected by 5-fold cross-validation on the training sample.

\textbf{Iterative Random Forest (iRF).}
\citeneeded{Basu et al 2018} adapted for coalition-game regression.
Binary coalition membership serves as features; the forest approximates
$v(S)$ through axis-aligned splits. Walsh coefficients are recovered
indirectly by applying the Walsh--Hadamard transform to the forest's
predictions over all $2^n$ coalitions. Hyperparameters (min leaf size,
max depth) selected by nested 3-fold CV on the training sample.

\textbf{KernelSHAP-IQ.}
\citeneeded{Fumagalli et al 2024} estimates $k$-Shapley Interaction
Indices ($k$-SII) \citeneeded{Grabisch 1997} up to a specified maximum
order via kernel-weighted regression. Smart sampling: the algorithm
chooses which coalitions to query. $k$-SII is a non-orthogonal basis;
reconstruction uses the plug-in predictor
$\hat{v}(S) = v(\emptyset) + \sum_{T \subseteq S, T \neq \emptyset} \phi_T$.

\textbf{SHAPIQ (Monte Carlo).}
\citeneeded{Muschalik et al 2024} estimates $k$-SII via Monte Carlo
sampling with stratified permutation-based estimators.

\textbf{SVARM-IQ.}
\citeneeded{Kolpaczki et al 2024} estimates $k$-SII via stratified
sampling with variance-aware allocation.


%% ====================================================================
\section{Evaluation Protocol}
\label{sec:protocol}
%% ====================================================================

\subsection{Paired-seed budget sweep}

All five methods are evaluated on identical pre-designated held-out sets
at each of seven budget levels (1\%, 2\%, 3\%, 5\%, 10\%, 20\%, 40\%
of $2^{15}$ coalitions). For each (circuit, budget, seed) triple:
%
\begin{enumerate}
\item A 10\% held-out set is drawn before any method runs.
\item LASSO-Walsh and iRF receive a random sample from the remaining
      pool. Shapiq methods choose their own coalitions via smart sampling
      (query overlap with the held-out set is tracked and logged).
\item All methods are scored on the \emph{same} held-out coalitions.
\end{enumerate}
%
Twenty seeds per budget level ($\leq 10\%$); ten seeds at higher budgets.

\subsection{Metrics}

\textbf{Reconstruction: held-out $R^2$.}
The primary reconstruction metric. Method-neutral because no method
optimizes held-out prediction directly.

\textbf{Reconstruction: per-order NMSE.}
Normalized mean squared error within each interaction order separately,
excluding order 0 (intercept) to avoid denominator domination.

\textbf{Detection: pairwise interaction AUROC.}
For each of the $\binom{15}{2} = 105$ head pairs, the true pairwise
Walsh coefficient magnitude $|w_{\{i,j\}}|$ defines a binary label
(above median = truly interacting). AUROC of the estimated
$|\hat{w}_{\{i,j\}}|$ predicting this label measures detection accuracy
independent of reconstruction quality.
%
\todo{Note: AUROC was added post-hoc as an exploratory metric. It was
not part of the pre-registered hypotheses. Predictions for AUROC were
generated by a blinded sub-agent before results were examined.}

\subsection{Pre-registration}

Hypotheses H1--H4 were frozen with SHA-256 hash
\texttt{2caafff\ldots} before the shapiq sweep ran. Version pins:
\texttt{shapiq==1.6.0}, \texttt{numpy==2.2.6},
\texttt{scipy==1.15.3}, \texttt{scikit-learn==1.6.1}.

\textbf{H1 (Estimation efficiency).}
All three shapiq approximators achieve $\geq 95\%$ of the exact $k$-SII
ceiling at 5\% budget.
\emph{Verdict:} \numplaceholder{Confirmed/Falsified.}

\textbf{H2 (Per-circuit ceiling behavior).}
Order-3 shapiq on \texttt{weight\_ioi} exceeds the order-2 ceiling
by $\geq 5$ $R^2$ points; canonical IOI shows $< 3$ points gain.
\emph{Verdict:} \numplaceholder{Confirmed/Falsified.}

\textbf{H3 (Method ordering).}
At $\geq 5\%$ budget, LASSO-Walsh $>$ iRF $>$ best shapiq (order 2)
on held-out $R^2$, with each gap $> 0.05$.
\emph{Verdict:} \numplaceholder{Confirmed/Falsified.}

\textbf{H4 (Cross-circuit consistency).}
Method ordering is consistent across all three circuits at $\geq 5\%$
budget.
\emph{Verdict:} \numplaceholder{Confirmed/Falsified.}

\todo{Add AUROC predictions from blind sub-agent here, with
contamination statement.}


%% ====================================================================
\section{Results}
\label{sec:results}
%% ====================================================================

\subsection{Reconstruction accuracy}

\begin{figure}[t]
\centering
\todo{Three-panel figure: R² vs budget for each circuit, 5 methods per panel.}
\caption{Held-out $R^2$ as a function of sample budget for five
interaction-detection methods across three ground-truth regimes. Error
bands show 95\% intervals over 20 seeds. LASSO-Walsh dominates at all
budgets on the near-additive and interactive circuits; the gap narrows
on the null circuit where all methods achieve high reconstruction by
default.}
\label{fig:r2-budget}
\end{figure}

\todo{Report the R² results. Key findings: LASSO dominates, shapiq
reaches ceiling efficiently, iRF has a floor from axis-aligned splits.}

\subsection{The ratio-metric trap}

\todo{The ``instructively wrong prediction'' story. Gap widens on easier
functions because LASSO's absolute error approaches zero while iRF hits
an invariant floor. Framed as why ratio metrics mislead---a
methodological contribution.}

\subsection{Estimation efficiency against ceilings}

\begin{figure}[t]
\centering
\todo{Efficiency plot: achieved R² / ceiling R² for shapiq methods vs budget.}
\caption{Estimation efficiency (achieved $R^2$ / exact $k$-SII ceiling)
for the three shapiq approximators. KernelSHAP-IQ reaches
\numplaceholder{98\%} of the ceiling at 5\% budget across all circuits,
confirming that the bottleneck is truncation order, not estimation
quality.}
\label{fig:efficiency}
\end{figure}

\todo{Report efficiency numbers. Key finding: near-perfect efficiency
at modest budgets across all approximators.}

\subsection{Pairwise interaction detection}

\begin{figure}[t]
\centering
\todo{AUROC vs budget, or AUROC bar chart at one representative budget.}
\caption{Pairwise interaction AUROC at \numplaceholder{5\%} budget.
LASSO-Walsh achieves near-perfect detection
(\numplaceholder{AUROC $\sim 1.0$}) because it directly estimates the
Walsh coefficients being scored. Shapiq methods, operating through
the $k$-SII basis, show a detection gap wider than their reconstruction
gap---basis misalignment degrades ranking more than aggregate prediction.}
\label{fig:auroc}
\end{figure}

\todo{Report AUROC results. Key finding: basis alignment matters more
for detection than reconstruction.}

\subsection{Per-order NMSE}

\begin{figure}[t]
\centering
\todo{Bar chart of per-order NMSE at one representative budget.}
\caption{Per-order NMSE at \numplaceholder{5\%} budget. \todo{Describe
which orders each method recovers well and poorly.}}
\label{fig:per-order-nmse}
\end{figure}


%% ====================================================================
\section{Metric Design: Three Failure Modes}
\label{sec:audit}
%% ====================================================================

% Promoted to a real section per outline. Each is a named failure mode.

Three metric-design bugs were identified and corrected during benchmark
development. Each represents a failure mode that recurs across the
interaction-detection and explainability literatures.

\subsection{Intercept domination}

\todo{NMSE including order-0 coefficient is dominated by the intercept
(97\%+ of energy). A method that returns only the mean achieves apparent
NMSE $< 0.03$. Fix: exclude order 0 from spectrum metrics.}

\subsection{Hyperparameter handicapping}

\todo{iRF with default min\_samples\_leaf=5 showed a puzzling inverted
curve. Nested CV on the training fold (not the held-out set) resolved
this: min\_leaf=1 was selected on 100\% of trials. The error floor that
remains after tuning is structural (axis-aligned splits vs.~global Walsh
rotation), not a hyperparameter artifact.}

\subsection{Ratio-metric explosion}

\todo{The ratio (LASSO NMSE / iRF NMSE) manufactured an apparent 89x
gap that disappeared under matched absolute scales. Diagnosis: LASSO's
numerator approached zero on near-additive circuits while iRF's floor
held constant. The ratio amplifies a denominator artifact into a
headline claim.}


%% ====================================================================
\section{Discussion}
\label{sec:discussion}
%% ====================================================================

\todo{Answer the genetics-facing question: which method should
practitioners trust? Under what regime?

Key points:
- LASSO-Walsh when interaction structure is sparse and basis is known
- shapiq methods when interpretable Shapley-style attribution is needed
  (trade detection for interpretability)
- iRF as a model-free baseline, but aware of the axis-alignment floor
- Basis choice dominates algorithm choice---this is the transferable insight}


%% ====================================================================
\section{Limitations}
\label{sec:limitations}
%% ====================================================================

\begin{itemize}
\item Three ground-truth regimes, not a systematic sweep of interaction
      density or sparsity.
\item Zero-ablation only; mean-ablation ground truth deferred.
\item Order $\leq 3$ truncation ceiling bounds all shapiq results;
      order $\geq 4$ interactions are not recovered.
\item Using an approximator (e.g., KernelSHAP-IQ) as ground truth at
      larger $n$ is explicitly \emph{not} validated by this work.
      Extrapolation would require independent verification that the
      approximator's error profile at $n = 15$ transfers to larger
      systems.
\item $n = 15$ is small relative to biological systems ($n \sim 10^3$
      to $10^6$ SNPs). The benchmark tests method \emph{accuracy},
      not scalability.
\end{itemize}


%% ====================================================================
\section{Conclusion}
\label{sec:conclusion}
%% ====================================================================

\todo{Basis alignment, not method sophistication, explains recovery
fidelity. Benchmark and pre-registration artifacts released.}


%% ====================================================================
%% References
%% ====================================================================

\bibliographystyle{plainnat}
% \bibliography{references}

\end{document}
